\section{Complete interface resolutions}\label{re:section}

The stopping construction uses a fixed face triangulation, whereas the
published loop limit uses leftmost interfaces. We must compare complete
collections, including the number and orientation of their macroscopic
loops: changing one face pairing can merge two loops or split one.

Throughout this section, locators for \cite{BH2019} refer to
arXiv version~3. We use its oriented-loop convention in Section~2.5: an
Ising circuit is a closed dual path, with no repeated edge, having plus
spins on its left and minus spins on its right.  It may visit a dual
vertex twice.  At a checkerboard face there are two noncrossing ways of
pairing the two incoming and two outgoing disagreement edges.  A
\emph{complete resolution} chooses a pairing at every such face and
uses every disagreement edge exactly once. Fixing a diagonal in each
face and taking the zero set of the affine spin interpolation on its
two triangles gives one such resolution, up to a displacement of order
$\delta$.

A loop is \emph{leftmost} when the adjacent plus spins along its left
side are connected successively by nearest-neighbor steps; this is the
strong-path convention of \cite[Section~2.5]{BH2019}. The canonical
leftmost collection consists of all loops satisfying that convention.

For oriented curves, let $d_{\mathrm{u}}$ denote uniform distance modulo
orientation-preserving parametrization; for loops also allow a change
of starting point.  Distances between loop collections are obtained by
matching loops and charging the diameters of unmatched loops, as in
\cite[Section~2.7]{BH2019}.

\begin{lemma}[Convergence of every complete resolution]
\label{if:resolution-convergence}
In a plus Jordan domain whose lattice boundary converges uniformly
modulo parametrization, every complete resolution of the spin
disagreement edges converges to the same oriented nested
$\mathrm{CLE}_3$ collection.  This assertion is uniform over the choice
of the resolution.  In particular it applies to a fixed choice of
diagonal in every square face.
\end{lemma}

\begin{proof}
Theorem~6 of \cite{BH2019} gives convergence of the leftmost loops in
the collection metric and the following simultaneous approximation.
For every $e>0$, with probability tending to one as $\delta\downarrow0$,
every Ising circuit $\ell$ of diameter greater than $e$ has a leftmost
circuit $\ell^{\mathrm L}$ such that
\[
 d_{\mathrm u}(\ell,\ell^{\mathrm L})\le e,
 \qquad
 \operatorname{diam} C<e
 \quad\text{for every component }C\text{ of }
 (\ell\cup\ell^{\mathrm L})\setminus(\ell\cap\ell^{\mathrm L}).
\]
The paragraph following that theorem explains the common-edge content:
the shared edges become dense in each macroscopic loop. In particular,
above a fixed positive diameter cutoff every circuit shares edges with
its matched leftmost circuit, with probability tending to one. These
assertions hold for all circuits on the same event, including circuits
selected after the spin configuration has been seen.

We explain why this one-sided approximation also supplies coverage and
the correct multiplicities for a complete resolution. From any vanishing
mesh sequence, first choose errors $e\downarrow0$ slowly enough that the
simultaneous events above have probabilities tending to one. Pass to a
subsequence and couple the leftmost collections to their limit so that
both their collection-distance error and the all-circuits error tend to
zero almost surely. Choose a diameter
cutoff which is not the diameter of a limit loop.  There are finitely
many limit loops above it.  They are simple and pairwise disjoint, and
distinct members of this finite family have positive mutual distance.

Any two complete resolutions are joined by changing the checkerboard
pairings one at a time.  Both possible pairings match incoming to
outgoing oriented edges.  Each change either merges two edge-simple
circuits or splits one into two edge-simple circuits.  Consequently all
intermediate circuits remain in the class covered by the all-circuits
assertion; no crossing circuit or repeated edge is introduced.

Suppose, along a sequence of finer meshes, that some such switch changes
the number of circuits matched near a fixed positive-diameter limit
loop.  Select one offending switch on each mesh.  The all-circuits
assertion holds for these selected circuits, even though the selection
depends on the configuration.  After passing to a subsequence, each
constituent either has vanishing diameter or converges, modulo
parametrization, to a simple loop in the limiting collection.  Two
vanishing constituents cannot merge to macroscopic size because they
meet at the switched vertex.  Adding a vanishing constituent to a large
one does not change the latter's parametrized limit.  If both
constituents are large, their meeting vertices force their limit loops
to meet.  The limits must therefore be the same simple loop.

The last alternative is impossible.  Choose a point strictly inside
that simple loop.  Uniform parametrized convergence preserves the
winding number about the point.  Each constituent has winding number
$1$ or $-1$, whereas the merged traversal has their sum, namely $0$ or
$\pm2$. The merged circuit is itself covered by the simultaneous
assertion. Its image approaches the same loop, so uniqueness of the
nearby limiting loop makes its \emph{oriented parametrized} limit that
loop as well. Its winding number must therefore be $\pm1$, a
contradiction. Splits
are the reverse operation.  It follows that no switch can change the
macroscopic count near any fixed limit loop, for all sufficiently fine
meshes on this coupled subsequence.  This argument selects an offending
switch if one exists and therefore does not accumulate an error over
the number of switches.

A specified leftmost circuit can be included in a complete resolution:
fix its pairings at every visited vertex and resolve the unused edges
arbitrarily. At a once-visited checkerboard vertex the unused incoming
and outgoing edges have the remaining noncrossing pairing. At a
twice-visited vertex the circuit already specifies both pairs. Thus its
oriented noncrossing traversal gives consistent local choices. In this
resolution a second circuit
matched to the same limit loop would have to share edges with the
specified leftmost circuit.  Uniqueness of the nearby macroscopic
leftmost loop and the common-edge consequence above give this
conclusion.  It contradicts use of every edge only once.  The count is
thus one in this resolution, and the switching argument makes it one
in every resolution.  The all-circuits assertion also excludes
unmatched positive-diameter circuits.  This proves convergence in the
collection metric.

Finally, colors are recovered by nesting from exterior plus.  Any
ancestor of a loop of diameter at least $a>0$ also has diameter at least
$a$.  There are only finitely many such ancestors.  Their nesting and
orientations are stable under convergence to disjoint simple loops,
so the colors needed at positive scales converge as well.
\end{proof}
